% Part: first-order-logic
% Chapter: sequent-calculus
% Section: provability-quantifiers

\documentclass[../../../include/open-logic-section]{subfiles}

\begin{document}

\olfileid{fol}{seq}{qpr}
\olsection{\usetoken{S}{derivability} en de kwantoren}

\begin{explain}
  Voor de volledigheidsstelling is bovendien nodig dat de regels van
  de sequentenrekening de in deze paragraaf bewezen feiten
  over~$\Proves$ opleveren.
\end{explain}

\begin{thm}
\ollabel{thm:strong-generalization} Als $c$ een constante is die niet
voorkomt in $\Gamma$ of $!A(x)$ en $\Gamma \Proves !A(c)$, dan geldt
$\Gamma \Proves \lforall[x][!A(x)]$.
\end{thm}

\begin{proof}
Laat $\pi_0$ een $\Log{LK}$-!!{derivation} zijn van $\Gamma_0 \Sequent
!A(c)$ voor een eindige $\Gamma_0 \subseteq \Gamma$. Door een
toepassing van $\RightR{\lforall}$ toe te voegen verkrijgen we
!!a{derivation} van $\Gamma_0 \Sequent \lforall[x][!A(x)]$. De
constante $c$ komt immers niet voor in $\Gamma$ of $!A(x)$, zodat aan
de eigenvariabelevoorwaarde is voldaan.
\end{proof}

\begin{prop}
\ollabel{prop:provability-quantifiers}
\begin{tagenumerate}{prvEx,prvAll}
\tagitem{prvEx}{$!A(t) \Proves \lexists[x][!A(x)]$.}{}
\tagitem{prvAll}{$\lforall[x][!A(x)] \Proves !A(t)$.}{}
\end{tagenumerate}
\end{prop}

\begin{proof}
\begin{tagenumerate}{prvEx,prvAll}
\tagitem{prvEx}{De sequent $!A(t) \Sequent \lexists[x][!A(x)]$ is
  !!{derivable}:
  \begin{prooftree}
    \Axiom$!A(t) \fCenter !A(t)$
    \RightLabel{\RightR{\lexists}}
    \UnaryInf$!A(t) \fCenter \lexists[x][!A(x)]$
    \end{prooftree}}{}
\tagitem{prvAll}{De sequent $\lforall[x][!A(x)] \Sequent !A(t)$ is
  !!{derivable}:
  \begin{prooftree}
    \Axiom$!A(t) \fCenter !A(t)$
    \RightLabel{\LeftR{\lforall}}
    \UnaryInf$\lforall[x][!A(x)] \fCenter !A(t)$
    \end{prooftree}}{}
\end{tagenumerate}
\end{proof}
\end{document}
